% 123-41(123쪽) — 곡선 y=-x^2+kx 와 x축 및 직선 x=4 로 둘러싸인 두 부분(색칠). 스캔 화질이 낮아 TikZ 로 다시 그림.
% 원문에 있는 것만: 라벨 O, x, y, y=-x^2+kx · x축 위 표시 k 와 4 · 색칠한 두 영역(0~k 는 x축 위, k~4 는 x축 아래).
% 눈금 숫자는 원문대로 k 와 4 뿐이다. k 의 자리는 원문 그림 비율 그대로이며 답이 되는 값은 적지 않는다.
\begin{tikzpicture}[x=0.40cm, y=0.26cm, line width=0.5pt]
  \draw[->, >=stealth, line width=0.7pt] (-1.1,0) -- (5.2,0) node[below=1pt] {$x$};
  \draw[->, >=stealth, line width=0.7pt] (0,-6.0) -- (0,4.2) node[left=1pt] {$y$};
  % 색칠한 두 영역
  \fill[yellow!45] plot[domain=0:2.7, samples=60, smooth] (\x, {-\x*\x+2.7*\x}) -- cycle;
  \fill[yellow!45] plot[domain=2.7:4, samples=40, smooth] (\x, {-\x*\x+2.7*\x}) -- (4,0) -- cycle;
  % 곡선
  \draw[line width=0.6pt, domain=-0.7:4.15, samples=140, smooth] plot (\x, {-\x*\x+2.7*\x});
  % 직선 x=4 (색칠한 부분의 오른쪽 경계)
  \draw[line width=0.5pt] (4,0) -- (4,-5.2);
  % 라벨(원문 자리 그대로)
  \node[below right=1pt and -3pt, font=\small] at (0,0) {$\pt{O}$};
  \node[below=2pt, font=\small] at (2.7,0) {$k$};
  \node[above left=0pt and -2pt, font=\small] at (4,0) {$4$};
  \node[anchor=west, font=\small] at (1.0,3.3) {$y=-x^2+kx$};
\end{tikzpicture}
